\scanpage{6}
où $C\in\unsure{\mathbb{S}|\mathrm{S}}(\mathbb{R})\simeq\mathbb{C}^{*}$ correspond à l'élément $\sqrt{-1}$ de $\mathbb{C}^{*}$.

[Ceci implique que $\varphi$ est alternée si $n$ impair, symétrique si $n$ pair]
\margin{\uline{NB} Soit $G'=\operatorname{Ker}\varepsilon$, \deleted{\unsure{On a} $\unsure{j(G_m)}$}, alors $G$ est \deleted{\unsure{un quotient de}} $\simeq G'\mathop{\times}\limits_{\mathbb{Z}/2\mathbb{Z}}G_{m\,\mathbb{Q}}$, où $\mathbb{Z}/2\mathbb{Z}\longrightarrow G'$ est induit par $i$. Donc $G$ (avec $i,\varepsilon$) se récupère en connaissant $G'$ et son point d'ordre 2 $i(-1)=\eta$. La donnée de \tikz[baseline=-.25ex,x=1ex,y=1ex]{\draw[line width=.5pt] (0,2.1) -- (0,.18) .. controls (0,-.12) and (.63,-.18) .. (.48,.12) .. controls (.40,.28) and (.18,.23) .. (0,.14);}}

\uline{Théorème} Pour $j$ donné, \unsure{relié|relé} à $i$ et $\varepsilon$ par condition a) et b), les conditions suivantes \inserted{(i)(ii)(i bis)(ii bis)} sont équivalentes et impliquent (iii) :

(i) Pour tout $V\in\operatorname{Ob}\operatorname{Rep}(G)$, \inserted{de degré $n$,} il existe une forme $\varphi:V\deleted{\unsure{\otimes}}\otimes V\longrightarrow\mathbb{Q}(-n)$ dans $\operatorname{Rep}(G)$ qui définit une polarisation (pour \deleted{\unsure{les}}la structure de Hodge \unsure{associées|associée|associés} à $j$ sur $V$, \unsure{$j$|$]$|$\}$} i.e. en fait : $C=j(i)$]

(i bis) Comme (i), mais \deleted{\unsure{avec}} pour des $V_{\alpha}$ donnés seulement \inserted{comme module sur $G'$}, \deleted{qui engendrent}\inserted{et} tels que \deleted{\unsure{les}} l'hom $G\longrightarrow\underline{\mathrm{Aut}}_{\mathbb{Q}}(V_{\alpha})$ soit injectif. De plus, on suppose $G'$ réductif i.e. $G$ réductif

(ii), \deleted{\unsure{(ii bis) Le centralisateur de}} Comme (i), (i bis), mais avec $G$ remplacé par $G_{\mathbb{R}}$, et $\operatorname{Rep}(G)$ par $\operatorname{Rep}(G_{\mathbb{R}})$.

(iii) \inserted{\deleted{\unsure{$G$}} $G'_{\mathbb{R}}$ est réductif, et si $K(\deleted{\unsure{j}})$ est} Le centralisateur de $j(C)$ \inserted{(noté encore $C$)} dans $G_{\mathbb{R}}^{\prime\,0}\deleted{\unsure{/i_{\mathbb{R}}(G_{m/\mathbb{R}}^{0})}}$\margin{\unsure{compos. de|composante} \deleted{\unsure{connexe}} topologique}, \inserted{l'image de $K(j)$ dans $G_{\mathbb{R}}^{0}/i_{\mathbb{R}}(\mathbb{R}^{*+})$} est un sous-groupe compact maximal, \deleted{\unsure{le}} i.e. le centralisateur de $C$ dans $G_{\mathbb{R}}^{\prime\,0}$ est un ss-groupe compact \deleted{\unsure{en groupe}}\inserted{maximal.}

% margin note written sideways
\margin{\uline{NB} (iii) équivaut aux conditions :

(iii bis) $G^{\mathrm{ad}}=G/\mathrm{Centre}$ avec le composé \deleted{\unsure{la}} $\mathbb{S}\mathbin{-}G\to G^{\inserted{\mathrm{ad}}}\deleted{\unsure{/\mathrm{Centre}|/G_{\mathrm{cen}}}}$ satisfait à (i)

(iii ter) $\exists$ un quotient \inserted{$G''$} de $G$ par un ss-gpe fini de $G'$, tel que le composé $\mathbb{S}\mathbin{-}G\rightharpoonupG''$ satisfasse à (i)}

\noindent\rule{\linewidth}{0.3pt}

\tikz[baseline=-.25ex,x=1ex,y=1ex]{\draw[line width=.5pt] (.34,2.5) -- (.28,.16) .. controls (-.02,.39) and (-.32,.02) .. (-.10,-.14) .. controls (.11,-.28) and (.29,-.03) .. (.28,.16);}\quad $j:\mathbb{S}\longrightarrow G_{\mathbb{R}}$ \deleted{équivaut} satisfaisant a) et b équivaut à la donnée d'un hom $j':\mathbb{U}\longrightarrow G_{\mathbb{R}}'$ satisfaisant $j'(-1)=\eta$
% archive sheet number: 5

